% !TEX program = xelatex
\documentclass[11pt,a4paper]{article}
\usepackage{fontspec}
\setmainfont{Calibri}
\setmonofont[Scale=0.88]{Consolas}
\newfontfamily\symbolfont{Segoe UI Symbol}
\newfontfamily\cjkfont{Microsoft YaHei}[Scale=0.9]
\usepackage{polyglossia}\setdefaultlanguage{polish}
\usepackage[table]{xcolor}
\usepackage{booktabs,longtable,array,geometry,enumitem,graphicx,fancyhdr,caption}
\usepackage[most]{tcolorbox}
\PassOptionsToPackage{hyphens}{url}
\usepackage[hidelinks,unicode]{hyperref}
\emergencystretch=2.5em
\geometry{a4paper,margin=2.1cm,top=2.3cm,bottom=2.3cm}
\setlength{\parindent}{0pt}\setlength{\parskip}{5pt}
\definecolor{apteal}{HTML}{0E7C86}
\definecolor{apteallight}{HTML}{E6F4F5}
\definecolor{apgrey}{HTML}{F3F4F6}
\definecolor{apamber}{HTML}{B45309}
\definecolor{apamberlight}{HTML}{FEF3C7}
\definecolor{apred}{HTML}{B91C1C}
\definecolor{apredlight}{HTML}{FEE2E2}
\definecolor{apgreen}{HTML}{15803D}
\definecolor{apgreenlight}{HTML}{DCFCE7}
\renewcommand{\arraystretch}{1.25}
\captionsetup{font=small,labelfont=bf}
\pagestyle{fancy}\fancyhf{}
\fancyfoot[L]{\small AMPERE POINT — rejestr usterek · Q11 z modułem Shelly X}
\fancyfoot[R]{\small\thepage}
\renewcommand{\headrulewidth}{0pt}
\newtcolorbox{uwaga}[1][Uwaga]{enhanced,breakable,colback=apamberlight,colframe=apamber,title={#1},fonttitle=\bfseries,boxrule=0.6pt,arc=2pt,left=6pt,right=6pt,top=4pt,bottom=4pt}
\newtcolorbox{info}[1][Jak to działa]{enhanced,breakable,colback=apteallight,colframe=apteal,title={#1},fonttitle=\bfseries,boxrule=0.6pt,arc=2pt,left=6pt,right=6pt,top=4pt,bottom=4pt}
\newtcolorbox{wazne}[1][Ważne]{enhanced,breakable,colback=apredlight,colframe=apred,title={#1},fonttitle=\bfseries,boxrule=0.6pt,arc=2pt,left=6pt,right=6pt,top=4pt,bottom=4pt}
\newtcolorbox{przyklad}[1][Przykład liczbowy]{enhanced,breakable,colback=apgreenlight,colframe=apgreen,title={#1},fonttitle=\bfseries,boxrule=0.6pt,arc=2pt,left=6pt,right=6pt,top=4pt,bottom=4pt}
\newtcolorbox{kodbox}{enhanced,breakable,colback=apgrey,colframe=apgrey,boxrule=0pt,arc=2pt,left=5pt,right=5pt,top=3pt,bottom=3pt,fontupper=\ttfamily\footnotesize}
\setlist{nosep,leftmargin=*}
\setcounter{secnumdepth}{2}
\begin{document}

{\LARGE\bfseries Q11 z modułem Shelly X: rejestr usterek\par}
\vspace{4pt}\textcolor{apteal}{\rule{\textwidth}{1.2pt}}

\emph{AMPERE POINT · prowadzony od 29 września 2026 · oprac. Dawid Cekała}


Wszystkie usterki i pułapki napotkane przy DevKicie (Q11 + Shelly X), także te już naprawione, żeby można było do nich wrócić. Każdy wpis mówi, co widać, skąd się bierze, jaki ma stan i gdzie jest opis. Numery się nie zmieniają; nowy wpis dostaje kolejny wolny numer (od S-37) i trafia do tabeli swojej grupy.


Stan: \leavevmode{\symbolfont\color{apred}\textbf{✘}} otwarte · \leavevmode{\symbolfont\color{apamber}\textbf{◐}} obejście albo częściowo · \leavevmode{\symbolfont\color{apgreen}\textbf{✔}} naprawione.


\section*{Nasz skrypt w module}

\begin{longtable}{>{\raggedright\arraybackslash}p{3.00cm}>{\raggedright\arraybackslash}p{3.43cm}>{\raggedright\arraybackslash}p{3.43cm}>{\raggedright\arraybackslash}p{1.10cm}>{\raggedright\arraybackslash}p{3.43cm}}
\toprule
\rowcolor{apteallight}\textbf{Nr} & \textbf{Objaw} & \textbf{Przyczyna} & \textbf{Stan} & \textbf{Dalej / opis} \\
\midrule\endhead
\bottomrule\endfoot
S-01 & Rozjazd pola i sterownika: pole pokazuje wartość, której sterownik nie ma. 12:02:44 tryb „do limitu energii” (sterownik został przy „od razu”), 13:44:27 wyłączenie ładowania z aplikacji (sterownik odpowiedział 2 razy „wł.”) & Gdy odpowiedź sterownika równa się wartości zapamiętanej w module, moduł loguje „same” i nie powiadamia skryptu. Próba sprawdzania po zapisie (v4.1) cofała też dobre wartości & \leavevmode{\symbolfont\color{apred}\textbf{✘}} & Znane odrzucenia (tryb przy limicie 0, okno o równych godzinach) blokuje v4.3. Na resztę potrzebna próba: co widzi skrypt w \texttt{getDatapoint().\allowbreak{}v} po odrzuceniu. Testy B3c, B1 w dokumencie 01 \\
S-02 & Spóźniona odpowiedź sterownika wracała do niego jako polecenie: limit energii 3 {\symbolfont→} 2 kWh (12:03:07), włącznik „wł.” bez polecenia (12:03:37) & Zmiana z HA i wpisanie starej odpowiedzi przez skrypt docierały do skryptu po sobie; skrypt brał drugą za nowe polecenie & \leavevmode{\symbolfont\color{apgreen}\textbf{✔}} v4.3 & Zmiany ze źródłem „sys” pomijane. Atrapa, test 11 \\
S-03 & Seria kliknięć w HA dawała serię zapisów: 24 polecenia w 9 s, sterownik przestawał odpowiadać & Każda zmiana pola szła od razu na łącze & \leavevmode{\symbolfont\color{apgreen}\textbf{✔}} v4.3 & Wysyłana ostatnia wartość po 1 s ciszy. Próba 13:27: 7 {\symbolfont→} 9 {\symbolfont→} 7 dało 1 zapis \\
S-04 & Przy każdym starcie modułu jedna zmiana pola „z zewnątrz” z wartością właśnie zgłoszoną przez sterownik & Źródło zdarzenia przy starcie nie jest „sys” (przyczyna nieustalona, startu nie widać w logu) & \leavevmode{\symbolfont\color{apgreen}\textbf{✔}} v4.3.1 & Zmiana na wartość, którą sterownik ma, nie idzie. Sprawdzić źródło, gdy będzie log startu (S-31) \\
S-05 & Okno godzin „1600” poszło jako 4 znaki ASCII, sterownik zignorował & Moduł wysyła tekst do punktu raw dosłownie, choć odczyt podaje szesnastkowo & \leavevmode{\symbolfont\color{apgreen}\textbf{✔}} v4.1 & 2 bajty przez \texttt{String.\allowbreak{}fromCharCode} \\
S-06 & Kontrola po zapisie cofnęła tryb „od razu” na „harmonogram” & Po własnym zapisie wartość widoczna dla skryptu nie była odświeżana & \leavevmode{\symbolfont\color{apgreen}\textbf{✔}} v4.2 & Kontrolę usunięto; wraca w S-01 \\
S-07 & Czas sesji skakał (770 {\symbolfont→} 815 {\symbolfont→} 1066 min) & Liczony z zegara modułu, który zmienia się po synchronizacji & \leavevmode{\symbolfont\color{apgreen}\textbf{✔}} v4 & Liczony z tyknięć co 10 s \\
S-08 & Wartości ułamkowe z HA: okno 0,01 h, limit 6,5 A & HA nie pilnuje kroku pola & \leavevmode{\symbolfont\color{apgreen}\textbf{✔}} v4.3 & Skrypt zaokrągla i wpisuje z powrotem \\
\end{longtable}

\section*{Moduł Shelly X i portal}

\begin{longtable}{>{\raggedright\arraybackslash}p{3.00cm}>{\raggedright\arraybackslash}p{3.81cm}>{\raggedright\arraybackslash}p{2.97cm}>{\raggedright\arraybackslash}p{0.81cm}>{\raggedright\arraybackslash}p{3.81cm}}
\toprule
\rowcolor{apteallight}\textbf{Nr} & \textbf{Objaw} & \textbf{Przyczyna} & \textbf{Stan} & \textbf{Dalej / opis} \\
\midrule\endhead
\bottomrule\endfoot
S-10 & Strona WWW modułu: kafelki faz „N/A A”, moc podpisana kW, choć jest w W & Szablon ekranu Shelly używa \texttt{displayValue} zamiast \texttt{value} & \leavevmode{\symbolfont\color{apred}\textbf{✘}} & Własny \texttt{web.\allowbreak{}svc.\allowbreak{}svelte} (panel) \\
S-11 & Strona WWW nie pokazuje 9 nowych pól & Szablon ekranu zna tylko 6 pól ładowarki & \leavevmode{\symbolfont\color{apred}\textbf{✘}} & Jak S-10 \\
S-12 & Akcje i warunki w manifeście: 2 i 3 zamiast 6 i 12 (kompilacje 3812–3815; 3811 miała 6/12) & Nieustalona; portal przy kompilacji bierze te z szablonu & \leavevmode{\symbolfont\color{apred}\textbf{✘}} & Sprawdzić przy B5 (webhooki) \\
S-13 & Moduł ponawia zapis po 0,3 s bez odpowiedzi („response to 07 timed out”). Ponowienia potrafią przestawić kolejność: 14:13:43–49 poszło 13, 12, 13, 13, 12 A; przez chwilę sterownik miał starsze 13 zamiast 12 & Stały czas oczekiwania w module, sterownik odpowiada wolniej (S-23); moduł ponawia stare polecenie po wysłaniu nowszego & \leavevmode{\symbolfont\color{apred}\textbf{✘}} & Dla Shelly: ponowienie nie powinno wyprzedzać nowszego zapisu tego samego punktu. U nas: odstępy między zapisami {\symbolfont≥} 3 s nie wystarczyły \\
S-14 & Zapis punktu, którego sterownik jeszcze nie zgłosił: „Setting unknown datapoint” & Moduł dodaje punkt dopiero po pierwszym meldunku & \leavevmode{\symbolfont\color{apamber}\textbf{◐}} & Skrypt podpina punkt później i zapisuje komunikat \\
S-15 & Otwarcie kreatora w portalu zakłada blokadę edycji; druga karta przechodzi w tryb tylko do odczytu & Blokada z sygnałem co 20 s & \leavevmode{\symbolfont\color{apamber}\textbf{◐}} & Nie otwierać kreatora, gdy Dawid edytuje; status \texttt{/\allowbreak{}products/\allowbreak{}3109/\allowbreak{}edit-\allowbreak{}lock/\allowbreak{}status} \\
S-16 & Portal przestawia wcięcia kilku linii skryptu & Edytor portalu & \leavevmode{\symbolfont\color{apamber}\textbf{◐}} & Porównywać treść bez białych znaków \\
S-17 & „Development Wi-Fi” zapisuje hasło Wi-Fi jawnie w tokenie produktu & Tak działa portal & \leavevmode{\symbolfont\color{apamber}\textbf{◐}} & Wgrywać samą usługę (Wi-Fi zostaje) \\
S-18 & Strona WWW i API bez hasła: każdy w sieci może sterować ładowarką & Hasło modułu wyłączone & \leavevmode{\symbolfont\color{apred}\textbf{✘}} & Test B11 \\
S-19 & Po restarcie bez internetu moduł nie ma czasu, a sterownik pyta o godzinę & Czas tylko z serwera w internecie & \leavevmode{\symbolfont\color{apred}\textbf{✘}} & Testy C; serwer czasu w sieci lokalnej \\
S-37 & Zmiana pola z harmonogramu modułu (i innych wywołań wewnątrz modułu) dociera do skryptu bez pola źródła; w logu \texttt{notifying for status:\allowbreak{} \{"value":\allowbreak{}6\}} (B4, 13:53:52) & Tak moduł zgłasza zmiany z wywołań lokalnych & \leavevmode{\symbolfont\color{apamber}\textbf{◐}} & Skrypt uznaje brak źródła za zmianę z zewnątrz, więc harmonogram działa. Prawdopodobnie to samo przy starcie (S-04) \\
S-50 & Po każdym restarcie modułu, około 40 s po nim, przychodzą webhooki „ładuje” i „start ładowania”, choć ładowanie trwa bez przerwy (30.09 15:01, 15:04, 15:07) & Pole stanu po starcie ma wartość domyślną „wolna”; meldunek sterownika zmienia je na „ładuje”, co wygląda jak nowy start & \leavevmode{\symbolfont\color{apred}\textbf{✘}} & Skrypt: po starcie nie wpisywać wartości domyślnej, tylko zostawić puste pole do pierwszego meldunku, albo webhooki z warunkiem na poprzedni stan; integracje liczące sesje muszą to odfiltrować \\
S-39 & Moduł zgłasza sterownikowi stan sieci 4 („połączony z chmurą”), choć chmura jest wyłączona, także bez internetu (co 30 s, bez zmian). Ikona Wi-Fi na wyświetlaczu co kilkadziesiąt sekund znika i wraca (30.09, poziom C) & Moduł wysyła stan Wi-Fi według połączenia z routerem, nie z chmurą & \leavevmode{\symbolfont\color{apred}\textbf{✘}} & \texttt{tmcu.\allowbreak{}overwriteWifiState} w skrypcie; test C5 \\
S-40 & Przy komplecie połączeń (HA, monitor, strona, wychodzące + 6) siódme nie wchodzi, a moduł zamyka połączenie Home Assistant & Limit połączeń w module (B9) & \leavevmode{\symbolfont\color{apred}\textbf{✘}} & Liczyć połączenia w projekcie mostu OCPP; HA łączy się ponownie sam \\
S-41 & Kopia konfiguracji niezaszyfrowana: hasło Wi-Fi i token chmury jawnie; nie zawiera plików usługi & Tak działa \texttt{Sys.\allowbreak{}CreateBackup} (B10) & \leavevmode{\symbolfont\color{apamber}\textbf{◐}} & Plik tylko w \texttt{DevKit\textbackslash{}kopie\_\allowbreak{}konfiguracji}; kopia nie odtworzy produktu na nowym module \\
S-42 & Kanał UDP (B12) steruje ładowarką bez hasła & Kanał bez uwierzytelnienia & \leavevmode{\symbolfont\color{apamber}\textbf{◐}} & Wyłączony po teście; sprawdzić przy B11, czy hasło go obejmuje \\
S-44 & Wartości pól zapamiętywanych (włącznik, limit prądu) moduł zapisuje do \texttt{storage.json} wprost, bez pliku tymczasowego: około 1,9 s po zmianie, najwyżej raz na około 3 s. Po zaniku przepadają zmiany z ostatnich około 3 s (30.09: 423 i 424 przepadły, wróciło 422). Częsta zmiana limitu zużywa pamięć: co 10 s granica około 2,7 roku, ciągłe zmiany poniżej roku & Tak moduł zapisuje pola z flagą „persisted” & \leavevmode{\symbolfont\color{apred}\textbf{✘}} & Dokument 04 (P4–P6). Zanik w samym zapisie nie trafiony, plik nieuszkodzony. Dla produktu: limit sterowany automatycznie bez flagi zapamiętywania; sterownik i tak pamięta swój limit \\
S-45 & Po zaniku zasilania, dopóki moduł nie dostanie godziny, wykonuje harmonogram według zapamiętanego, nieaktualnego zegara (30.09: start z zegarem 29.09 13:26:55, zadania co 5 s odpalały „o 13:27:00, 13:27:35…”), choć raportuje godzinę jako nieznaną & Moduł bez baterii zegara startuje od zapamiętanej chwili i nie wstrzymuje harmonogramu & \leavevmode{\symbolfont\color{apred}\textbf{✘}} & Groźne na poziomie C: zadanie „22:00 zacznij ładować” odpali o złej porze. Serwer czasu w sieci lokalnej skraca okno do około 25 s; zgłosić do Shelly \\
\end{longtable}

\section*{Sterownik Q11}

\begin{longtable}{>{\raggedright\arraybackslash}p{3.00cm}>{\raggedright\arraybackslash}p{4.13cm}>{\raggedright\arraybackslash}p{2.33cm}>{\raggedright\arraybackslash}p{0.81cm}>{\raggedright\arraybackslash}p{4.13cm}}
\toprule
\rowcolor{apteallight}\textbf{Nr} & \textbf{Objaw} & \textbf{Przyczyna} & \textbf{Stan} & \textbf{Dalej / opis} \\
\midrule\endhead
\bottomrule\endfoot
S-20 & Po zapisie limitu energii sterownik przez kilka sekund powtarza meldunek limitu i trybu, około 10 par na sekundę (12:02:00–03 {\symbolfont≈} 45 par), i nie odpowiada modułowi (6 razy „Heartbeat timed out”, każde ustąpiło po 1–5 s) & Oprogramowanie sterownika & \leavevmode{\symbolfont\color{apamber}\textbf{◐}} & Skrypt ogranicza zapisy (S-03). Pytanie do fabryki \\
S-21 & Tryb pracy zmienia się sam: limit energii > 0 {\symbolfont→} „do limitu energii”, okno {\symbolfont→} „harmonogram”, „od razu” zeruje limit; „do limitu energii” przy limicie 0 odrzucony & Tak działa sterownik (29.09) & \leavevmode{\symbolfont\color{apamber}\textbf{◐}} & Skrypt nie wysyła trybu, który zostanie odrzucony. \texttt{DANE\_\allowbreak{}DIAGNOSTYCZNE.\allowbreak{}md} \\
S-22 & Okno o równej godzinie startu i końca (0–0, 1–1) odrzucane & Sterownik & \leavevmode{\symbolfont\color{apamber}\textbf{◐}} & Skrypt go nie wysyła; stan fabryczny 0–0 nie do przywrócenia z modułu \\
S-23 & Odpowiedź po 0,3–3 s i często najpierw ze starą wartością (limit 14 {\symbolfont→} zapis 9 {\symbolfont→} meldunki 14, 9) & Sterownik & \leavevmode{\symbolfont\color{apamber}\textbf{◐}} & Obsłużone w v4.3 \\
S-24 & Sterownik nie przyjmuje wyłączenia ładowania: 29.09 13:44:27 (z aplikacji), 30.09 10:59:56 i 11:06:50 (przez API). Restart samego modułu pokazał, że sterownik dalej ma „wł.”. 29.09 o 12:03 takie samo wyłączenie przyjął. Sam ustawia włącznik: „wł.” po włączeniu zasilania, „wył.” po zakończeniu sesji & Nieustalona; może zależeć od trybu albo stanu sterownika & \leavevmode{\symbolfont\color{apred}\textbf{✘}} & Porównać warunki z 12:03 (tryb „do limitu energii”) z późniejszymi; pytanie do fabryki \\
S-25 & Mapa usterek (punkt 10) ma 2 bajty, definicja 17 kodów; bit 16 „przegrzanie” się nie mieści & Sterownik albo definicja & \leavevmode{\symbolfont\color{apred}\textbf{✘}} & Pytanie do fabryki \\
S-26 & Limit prądu w definicji Tuya 6–32 A, Q11 to 16 A & Definicja produktu & \leavevmode{\symbolfont\color{apamber}\textbf{◐}} & Skrypt przyjmuje 6–16 A \\
S-46 & Po restarcie modułu łącze ze sterownikiem utknęło w stanie „Fetching Config” na ponad 3 min (30.09 12:24–12:27): sterownik odpowiadał na sygnał życia i pytał o godzinę, ale moduł nie dokończył powitania i nie ponawiał pytania; brak danych i sterowania. Pomógł kolejny restart. Zdarzyło się przy pierwszym restarcie modułu po zaniku zasilania & Warstwa łącza w module (Shelly) nie ponawia pytania o konfigurację & \leavevmode{\symbolfont\color{apred}\textbf{✘}} & Zgłosić do Shelly. U nas: w skrypcie strażnik — gdy przez 2 min po starcie brak meldunków sterownika, komunikat w stanie usługi i ewentualnie restart modułu \\
S-47 & Przy zaniku zasilania sterownik traci ustawienie jednej sesji na czas: okno godzin z trybem „harmonogram” (wraca „od razu”) oraz opóźnienie i czas trwania z menu (sprawdzone 2 razy 30.09); limit prądu pamięta. Klient, który ustawił opóźnienie na tańszą taryfę, po mrugnięciu prądu ładuje od razu & Sterownik nie zapisuje ustawień sesji & \leavevmode{\symbolfont\color{apred}\textbf{✘}} & Pytanie do fabryki. Propozycja (niesprawdzona): moduł zapamiętuje ustawienie sesji i odtwarza je po starcie. To nie jest harmonogram — harmonogram Tuya ustawia się w aplikacji \\
S-49 & Opóźnienie i czas trwania ustawiane w menu, gdy ładowarka jest w stanie C, pokazują się na ekranie z innymi wartościami niż wybrane (Dawid, 30.09). Występuje też na ładowarce z modułem Tuya & Sterownik (nie moduł) & \leavevmode{\symbolfont\color{apred}\textbf{✘}} & Powtarzalne niezależnie od ustawionych godzin; tak samo z Shelly i z WBR3, więc nie różnicuje modułów. Odłożone jako osobny, większy temat do fabryki (Dawid, 30.09) \\
S-48 & Opóźnienie i czas trwania ładowania ustawione w menu ładowarki (30.09 ok. 12:35) są dla modułu niewidoczne: sterownik nic nie wysłał, a przy odpytaniu podał okno 0–0 i tryb „od razu”. Widać tylko skutek: z testerem stan „czeka” mimo sygnału auta „chce ładować” (12:41) & Funkcja menu nie ma odpowiednika w zgłaszanych punktach; kandydat: punkt 33 \texttt{mode\_set} (surowy, do 128 B), nigdy nie zgłoszony & \leavevmode{\symbolfont\color{apred}\textbf{✘}} & Z testerem: czy opóźnienie z menu działa i czy moduł widzi „czeka”. Pytanie do fabryki o punkt 33 \\
S-27 & Punkt 33 nigdy nie przyszedł; licznik (1) tylko w stanie „ładuje”, co 90 s. Wersję (23) „V1” sterownik zgłasza po włączeniu zasilania (30.09: 3 razy). Po włączeniu zasilania nie zgłasza okna godzin (19), także przy późniejszym restarcie samego modułu, więc pola okna zostają 0–0 & Sterownik & \leavevmode{\symbolfont\color{apamber}\textbf{◐}} & Zestawienie DP; okno: sprawdzić, czy sterownik je kasuje przy zaniku (ustawić okno, zanik, odczyt z menu) \\
S-38 & Zapis trybu „od razu” (14:01:43) bez żadnej odpowiedzi, a sterownik go przyjął (odpytanie przy starcie o 14:04: tryb 0). O 12:02:38 ten sam zapis potwierdził & Sterownik nie zawsze potwierdza zapis & \leavevmode{\symbolfont\color{apred}\textbf{✘}} & Razem z S-01: brak potwierdzenia nie znaczy odmowy \\
\end{longtable}

\section*{Aplikacja Shelly, Home Assistant, narzędzia}

\begin{longtable}{>{\raggedright\arraybackslash}p{3.00cm}>{\raggedright\arraybackslash}p{4.04cm}>{\raggedright\arraybackslash}p{2.99cm}>{\raggedright\arraybackslash}p{0.81cm}>{\raggedright\arraybackslash}p{3.56cm}}
\toprule
\rowcolor{apteallight}\textbf{Nr} & \textbf{Objaw} & \textbf{Przyczyna} & \textbf{Stan} & \textbf{Dalej / opis} \\
\midrule\endhead
\bottomrule\endfoot
S-30 & Aplikacja Shelly bez chmury „offline”; w sieci lokalnej tylko sprawdza moduł & Aplikacja steruje wyłącznie przez chmurę (próba 13:43–13:45: z chmurą działa od razu) & \leavevmode{\symbolfont\color{apred}\textbf{✘}} & Ograniczenie Shelly; bez chmury strona WWW i HA. Test B1 \\
S-31 & Startu modułu nie widać w logu & Monitor łączy się 20–45 s po starcie; log przez MQTT też za późno i zapełnia się własnymi wpisami & \leavevmode{\symbolfont\color{apamber}\textbf{◐}} & Log UDP (odbiornik w kontenerze, port 8514) pokazuje start od 23. sekundy, czyli od połączenia z Wi-Fi; wcześniejszych sekund nie widać żadną drogą sieciową \\
S-43 & Zegar laptopa spóźnia się o 8,3 s (serwer czasu w kontenerze), więc czasy w logach po stronie laptopa (webhooki, WebSocket, zdarzenia) są o 8 s wcześniejsze niż w logu modułu & Zegar Windows & \leavevmode{\symbolfont\color{apamber}\textbf{◐}} & Przy porównywaniu logów przesuwać o 8 s albo zsynchronizować zegar laptopa (decyzja Dawida) \\
S-32 & Aplikacja rysuje ogólny ekran ładowarki, nie ekran TopAC & Ekran zależy od typu produktu w aplikacji & \leavevmode{\symbolfont\color{apred}\textbf{✘}} & Do ustalenia z Shelly \\
S-33 & Oficjalna integracja HA pokazuje tylko fazy (19 encji) & Encje tylko dla kluczy z listy w kodzie; limit i włącznik tylko dla TopAC (EVE01) & \leavevmode{\symbolfont\color{apamber}\textbf{◐}} & Łatka \texttt{DevKit\textbackslash{}ha\_\allowbreak{}integracja} (33 encje); docelowo zmiana w integracji \\
S-34 & Integracja HA wywraca się na polu wyboru bez napisów opcji & Błąd w kodzie integracji (\texttt{titles}) & \leavevmode{\symbolfont\color{apamber}\textbf{◐}} & Naprawione w łatce; zgłosić do HA \\
S-35 & Łatana kopia integracji HA po aktualizacji HA może nie wstać & Kopia z wersji 2026.6.4 & \leavevmode{\symbolfont\color{apamber}\textbf{◐}} & \texttt{ha\_\allowbreak{}integracja\textbackslash{}CZYTAJ.\allowbreak{}pdf} \\
S-36 & Polskie tłumaczenie aplikacji z błędami („Oskarżony”) & Aplikacja Shelly & \leavevmode{\symbolfont\color{apred}\textbf{✘}} & Analiza GAP \\
\end{longtable}

\end{document}